%
% 11-polynome.tex -- Potenzreihen von Polynomen
%
% (c) 2026 Prof Dr Andreas Müller
%
\subsection{Potenzreihen von Polynomen
\label{buch:analytisch:potenzreihen:subsection:polynome}}
Das Polynom
\[
p(z)
=
p_0
+ 
p_1z
+
p_2z^2
+
\dots
+
p_{n-1}z^{n-1}
+
p_nz^n
=
\sum_{k=0}^n p_kz^k
\]
ist eine Potenzreihe, die bereits nach endlich vielen Termen
abbricht.
\index{Polynom}%
Wir möchten den Zusammenhang zwischen den Koeffizienten $p_k$
und der Taylor-Reihe von $p(z)$ herstellen.
\index{Taylor-Reihe}%
\index{Potenzreihe}
Dazu benötigen wir eine Formel für die $k$-te Ableitung einer
Potenzfunktion, die wir im folgenden Lemma aus der Produktregel
ableiten.

\begin{lemma}
\label{buch:analytisch:potenzreihen:lemma:ableitung}
Die $k$-te Ableitung von $z^n$ ist
\begin{equation}
\frac{d^k}{dz^k}
z^n
=
\underbrace{
n(n-1)\cdots(n-k+1)
}_{\textstyle\text{$k$ Faktoren}}
z^{n-k}.
\label{buch:analytisch:potenzreihen:eqn:ableitung}
\end{equation}
\end{lemma}

\begin{proof}
Wir beweisen die Formel
\eqref{buch:analytisch:potenzreihen:eqn:ableitung}
durch vollständige Induktion.
Für $k=0$ wird gar nicht abgeleitet, die Formel gilt trivialerweise.

Für den Induktionsschritt nehmen wir jetzt an, dass die Formel für die
$k$-te Ableitung bereits bewiesen ist, und zeigen kann, dass sie auch
für die $(k+1)$-te Ableitung gilt.
Dazu berechnen wir
\begin{align*}
\frac{d^{k+1}}{dz^{k+1}} z^n
&=
\frac{d}{dz}
\frac{d^k}{dz^k}
z^n
\\
&=
\frac{d}{dz}
n(n-1)\dots(n-k+1) z^{n-k}
\\
&=
n(n-1)\dots(n-k+1)(n-k) z^{n-k-1}
\\
&=
n(n-1)\dots(n-k+1)(n-(k+1)+1) z^{n-(k+1)}.
\end{align*}
Dies ist die Formel für die $(k+1)$-te Ableitung.
Damit ist der Induktionsschritt gelungen und damit die Formel
bewiesen.
\end{proof}

Das Polynom $p(z)$ kann auch als Potenzreihe im Punkt $z_0$
geschrieben werden.
Dies kann auf zwei Arten erreicht werden.
Einerseits gilt die Taylor-Reihe
\begin{align*}
p(z)
&=
\sum_{u=0}^n
\frac{1}{u!}
\frac{d^up}{dz^u}(z_0)
(z-z_0)^u.
\intertext{Andererseits kann man eine Potenzreihenentwicklung durch eine
algebraische Umformung erhalten.
Zunächst wendet man auf}
p(z)
&=
p\bigl((z-z_0) + z_0\bigr)
=
\sum_{k=0}^n p_k \bigl((z-z_0)+z_0\bigr)^k
\intertext{die binomische Formel an und erhält}
&=
\sum_{k=0}^n
p_k
\sum_{u+v=k}
\binom{k}{u}
(z-z_0)^u
z_0^v.
\intertext{Durch Umstellung der Summe und Expansion des
Binomialkoeffizienten kann man dies in dir Form}
&=
\sum_{u=0}^n
\Biggl(
\sum_{k=u}^n
p_k
\binom{k}{u}
z_0^{k-u}
\Biggr)
(z-z_0)^u
\\
&=
\sum_{u=0}^n
\Biggl(
\sum_{k=u}^n
p_k
\frac{k(k-1)\cdots(k-u+1)}{u!}
z_0^{k-u}
\Biggr)
(z-z_0)^u
\intertext{bringen.
Der Nenner $u!$ kann aus der inneren Summe herausgezogen werden und
die verbleibenden Terme können mit Hilfe der Formel
\eqref{buch:analytisch:potenzreihen:eqn:ableitung}
in die Form}
&=
\sum_{u=0}^n
\frac{1}{u!}
\Biggl(
\sum_{k=u}^n
p_k
k(k-1)\cdots(k-u+1)
z_0^{k-u}
\Biggr)
(z-z_0)^u
\\
&=
\sum_{u=0}^n
\frac{1}{u!}
\Biggl(
\sum_{k=u}^n
p_k
\frac{d^u}{dz_0^u}z_0^k
\Biggr)
(z-z_0)^u
\intertext{gebracht werden.
Die Ableitung kann aus der inneren Summe herausgezogen werden,
so dass sich die Ableitung}
&=
\sum_{u=0}^n
\frac{1}{u!}
\frac{d^u}{dz_0^u}
\Biggl(
\sum_{k=0}^n
p_k
z_0^k
\Biggr)
(z-z_0)^u
\\
&=
\sum_{u=0}^n
\frac{1}{u!}
\frac{d^up}{dz_0^u}
(z_0)
(z-z_0)^u
\end{align*}
des Polnoms $p(z_0)$ an der Stelle $z_0$ ergibt.
Die algebraisch gefundene Potenzreihe stimmt also automatisch
mit der Taylor-Reihe überein.
